% !TeX program = lualatex
% !TeX encoding = utf8
% !TeX spellcheck = uk_UA
% !BIB program = biber

\documentclass{article}

\usepackage[%
	a4paper,%
	footskip=1cm,%
	headsep=0.3cm,%
	top=2cm, %поле зверху
	bottom=2cm, %поле знизу
	left=2cm, %поле ліворуч
	right=2cm, %поле праворуч
    ]{geometry}

\usepackage[fontsize=14pt]{fontsize}
\usepackage{fontsetup}
\usepackage[english, russian, ukrainian]{babel}

\usepackage{amsmath}
\usepackage[most, many]{tcolorbox}

% Визначаємо власні команди для зручності
\newcommand{\pd}[2]{\frac{\partial #1}{\partial #2}}
\newcommand{\pdd}[2]{\frac{d^2 #1}{d #2^2}}

% --- "Старовинне" оформлення боксів-приміток: пергамент, сепія, подвійна рамка ---
\definecolor{parchment}{RGB}{240,224,188}
\definecolor{sepiaframe}{RGB}{92,58,26}
\definecolor{sepiaink}{RGB}{74,44,18}

\tcbset{
  vintagebox/.style={
    enhanced,
    breakable,
    colback=parchment,
    colupper=sepiaink,
    colframe=sepiaframe,
    boxrule=0.7pt,
    arc=0pt,
    left=7mm, right=7mm, top=6mm, bottom=5mm,
    borderline={0.4pt}{-3.2pt}{sepiaframe},        % внутрішня тонка лінія -- подвійна рамка, як у старій книзі
    drop fuzzy shadow=black!50!white,              % легка тінь -- наче фоліант лежить на столі
    fonttitle=\scshape\normalsize,
    coltitle=sepiaink,
    colbacktitle=parchment,
    center title,
    top=8mm,
    attach boxed title to top center={yshift=-3.2mm},
    boxed title style={
        colback=parchment, colframe=sepiaframe, boxrule=0.5pt,
        arc=0pt, boxsep=1pt, left=2mm, right=2mm,
    },
    fontupper=\itshape\footnotesize,
  }
}

% Оформлення для приміток-пояснень (додано перекладачем для прояснення матеріалу)
\newtcolorbox{note}[1][]{%
    vintagebox,
    title={\upshape\rule{4mm}{0.3pt}~~Пояснення~~\rule{4mm}{0.3pt}},
    #1
}

% Початок документа
\begin{document}

\section*{Розв'язання рівняння Шрединґера для задачі Кеплера}

\begin{note}
\textbf{Про що йдеться.} Задача Кеплера в квантовій механіці --- це рух електрона в кулонівському полі ядра із зарядом $Ze$ (тобто водневоподібний атом: водень при $Z=1$, іонізований гелій при $Z=2$ тощо). Її класичним аналогом є рух планети навколо Сонця за законом тяжіння Ньютона (звідси й назва --- задача Кеплера), оскільки і кулонівська, і гравітаційна сили спадають як $1/r^2$. Мета цього розділу --- розв'язати рівняння Шрединґера для такої системи і показати, що енергія електрона може набувати лише дискретного набору значень $E_n$ --- саме тих, які раніше (в старій квантовій теорії) отримав Бор із напівкласичних міркувань.

Нагадаємо позначення: $\hbar = h/2\pi$ --- стала Планка, поділена на $2\pi$; $m$ --- маса електрона; $e$ --- елементарний заряд; $\Delta$ --- оператор Лапласа.
\end{note}

Рівняння Шрединґера для задачі Кеплера має вигляд

\begin{equation}
\Delta \psi + \frac{8\pi^2 m}{h^2} \left( E + \frac{e^2Z}{r} \right) \psi = 0.
\end{equation}

Вводячи сферичні координати \( r, \theta, \varphi \), приведемо його до вигляду

\begin{equation}
\left\{ \frac{\partial^2}{\partial r^2} + \frac{2}{r} \frac{\partial}{\partial r} + \frac{1}{r^2} \frac{\partial^2}{\partial \theta^2} + \frac{\cot \theta}{r^2} \frac{\partial}{\partial \theta} + \frac{1}{r^2 \sin^2 \theta} \frac{\partial^2}{\partial \varphi^2} \right\} \psi +
\frac{8\pi^2 m}{h^2} \left( E + \frac{e^2Z}{r} \right) \psi = 0.
\end{equation}

У цьому рівнянні змінні знову розділяються; якщо покласти

\begin{equation}
\psi = R(r) \Theta(\theta) \Phi(\varphi),
\end{equation}

\begin{tcolorbox}[vintagebox, title={\upshape\rule{4mm}{0.3pt}~~Детально: як саме відбувається розщеплення~~\rule{4mm}{0.3pt}}]

\textbf{Крок 1. Підставляємо $\psi = R\Theta\Phi$.}
Кожен доданок оператора діє лише на «свою» змінну ($r$, $\theta$ або $\varphi$), тому

\begin{equation*}
\Theta\Phi\left(R'' + \frac{2}{r}R'\right)
+ \frac{R\Phi}{r^2}\left(\Theta'' + \cot\theta\,\Theta'\right)
+ \frac{R\Theta}{r^2\sin^2\theta}\Phi''
+ \frac{8\pi^2 m}{h^2}\!\left(E+\frac{e^2Z}{r}\right)\! R\Theta\Phi = 0,
\end{equation*}
де штрих означає похідну за відповідною змінною ($R'=dR/dr$ тощо).

\textbf{Крок 2. Ділимо на $R\Theta\Phi$.}
Це стандартний прийом методу розділення змінних: ділення на добуток самих функцій прибирає їх із чисельника і залишає лише логарифмічні похідні:

\begin{equation*}
\frac{R''}{R} + \frac{2}{r}\frac{R'}{R}
+ \frac{1}{r^2}\left(\frac{\Theta''}{\Theta} + \cot\theta\,\frac{\Theta'}{\Theta}\right)
+ \frac{1}{r^2\sin^2\theta}\frac{\Phi''}{\Phi}
+ \frac{8\pi^2 m}{h^2}\!\left(E+\frac{e^2Z}{r}\right) = 0.
\end{equation*}

\textbf{Крок 3. Відокремлюємо $r$ від $(\theta,\varphi)$.}
Помножимо все на $r^2$ і перенесемо доданки, що залежать від $r$, в один бік:

\begin{equation*}
\underbrace{r^2\left[\frac{R''}{R} + \frac{2}{r}\frac{R'}{R} + \frac{8\pi^2 m}{h^2}\!\left(E+\frac{e^2Z}{r}\right)\right]}_{\text{залежить лише від } r}
= -\underbrace{\left[\frac{\Theta''}{\Theta} + \cot\theta\,\frac{\Theta'}{\Theta} + \frac{1}{\sin^2\theta}\frac{\Phi''}{\Phi}\right]}_{\text{залежить лише від } \theta,\varphi}.
\end{equation*}

Ліва частина -- функція тільки від $r$, права -- тільки від $\theta$ і $\varphi$. А оскільки $r$, $\theta$, $\varphi$ -- незалежні змінні, рівність можлива лише тоді, коли обидві частини дорівнюють одній і тій самій сталій. Позначимо цю сталу (константу розділення) через $\lambda$:

\begin{equation*}
r^2\left[\frac{R''}{R} + \frac{2}{r}\frac{R'}{R} + \frac{8\pi^2 m}{h^2}\!\left(E+\frac{e^2Z}{r}\right)\right] = \lambda,
\qquad
\frac{\Theta''}{\Theta} + \cot\theta\,\frac{\Theta'}{\Theta} + \frac{1}{\sin^2\theta}\frac{\Phi''}{\Phi} = -\lambda .
\end{equation*}

Перше зі співвідношень одразу дає радіальне рівняння (4):
\begin{equation*}
\left\{ \frac{d^2}{dr^2} + \frac{2}{r} \frac{d}{dr} + \frac{8\pi^2 m}{h^2} \left( E + \frac{e^2Z}{r} \right) - \frac{\lambda}{r^2} \right\} R = 0.
\end{equation*}

\textbf{Крок 4. Відокремлюємо $\theta$ від $\varphi$.}
Залишається розщепити друге співвідношення. Помножимо його на $\sin^2\theta$:

\begin{equation*}
\sin^2\theta\,\frac{\Theta''}{\Theta} + \sin\theta\cos\theta\,\frac{\Theta'}{\Theta} + \lambda\sin^2\theta = -\frac{\Phi''}{\Phi}.
\end{equation*}

Знову: ліва частина залежить тільки від $\theta$, права -- тільки від $\varphi$, отже обидві дорівнюють одній сталій, яку позначимо $m^2$:

\begin{equation*}
-\frac{\Phi''}{\Phi} = m^2
\quad\Longrightarrow\quad
\Phi'' + m^2 \Phi = 0
\end{equation*}
-- це і є рівняння (6);

\begin{equation*}
\sin^2\theta\,\frac{\Theta''}{\Theta} + \sin\theta\cos\theta\,\frac{\Theta'}{\Theta} + \lambda\sin^2\theta = m^2.
\end{equation*}

Ділимо останнє на $\sin^2\theta$:
\begin{equation*}
\frac{\Theta''}{\Theta} + \cot\theta\,\frac{\Theta'}{\Theta} + \lambda - \frac{m^2}{\sin^2\theta} = 0
\quad\Longrightarrow\quad
\left\{ \frac{d^2}{d\theta^2} + \cot \theta \frac{d}{d\theta} + \lambda - \frac{m^2}{\sin^2 \theta} \right\} \Theta = 0,
\end{equation*}
-- це рівняння (5).

\textbf{Підсумок.} Одне рівняння в частинних похідних для $\psi(r,\theta,\varphi)$ звелося до трьох звичайних диференціальних рівнянь -- по одному для кожної змінної, -- пов'язаних між собою двома сталими розділення: $\lambda$ (що з'явиться далі як $l(l+1)$) та $m^2$ (де $m$ -- магнітне квантове число). Це можливо саме тому, що змінні в лапласіані сферичних координат і в кулонівському потенціалі \emph{розділяються}: кожен доданок рівняння залежить не більш ніж від однієї змінної після відповідного групування.
\end{tcolorbox}

то рівняння можна розщепити на три рівняння:

\begin{equation}
\left\{ \frac{d^2}{dr^2} + \frac{2}{r} \frac{d}{dr} + \frac{8\pi^2 m}{h^2} \left( E + \frac{e^2Z}{r} \right) - \frac{\lambda}{r^2} \right\} R = 0,
\end{equation}

\begin{equation}
\left\{ \frac{d^2}{d\theta^2} + \cot \theta \frac{d}{d\theta} + \lambda - \frac{m^2}{\sin^2 \theta} \right\} \Theta = 0,
\end{equation}

\begin{equation}
\frac{d^2 \Phi}{d\varphi^2} + m^2 \Phi = 0,
\end{equation}
%
де \( m \) (як і раніше) та \( \lambda \) --- параметри розділення. Розв'язок третього рівняння вже відомий:

\begin{equation}
\Phi = \begin{bmatrix}
\cos(m\varphi)\\
\sin(m\varphi)
\end{bmatrix}
, \text{або} \quad \Phi = e^{im\varphi},
\end{equation}
%
де \( m \) має бути цілим числом, інакше розв'язок \( \Phi \) виявиться неоднозначним.

\begin{note}
Тут \(m\) --- це параметр розділення змінних у рівнянні для $\Phi(\varphi)$, а не маса електрона (в тому ж рівнянні вище маса позначена тією самою літерою $m$ --- це незручність позначень, успадкована з оригіналу). Надалі $m$ означатиме саме \emph{магнітне квантове число}: воно ціле, оскільки хвильова функція $\Phi(\varphi)=e^{im\varphi}$ має бути однозначною при повороті на $2\pi$, тобто $\Phi(\varphi+2\pi)=\Phi(\varphi)$.
\end{note}

Друге рівняння --- це рівняння, що визначає кульові функції (приєднані поліноми Лежандра) $P_l^m(\cos \theta)$, де $\lambda$ набуває значень $l(l+1)$, а $|m| \leq l$. Для інших значень рівняння не має скінченних однозначних розв'язків. Доведемо це в загальному вигляді, для чого об'єднаємо залежність від $\theta$ та від $\phi$ в узагальнену кульову функцію $Y_l(\theta, \phi)$. Якщо для стислості ввести позначення

\begin{equation}
\Lambda = \frac{1}{\sin \theta} \frac{\partial}{\partial \theta} \sin \theta \frac{\partial}{\partial \theta} + \frac{1}{\sin^2 \theta} \frac{\partial^2}{\partial \phi^2},
\end{equation}
%
так що

\begin{equation}
\Delta = \frac{\partial^2}{\partial r^2} + \frac{2}{r} \frac{\partial}{\partial r} + \frac{\Lambda}{r^2},
\end{equation}
%
то стане зрозуміло, що функція $Y_l(\theta, \phi)$ повинна задовольняти диференціальному рівнянню

\begin{equation}
\Lambda Y_l + \lambda Y_l = 0.
\end{equation}

Загальний розв'язок цього рівняння знаходиться так. Розглянемо однорідні поліноми $U_l$ степеня $l$ за $x$, $y$, $z$, що задовольняють рівнянню Лапласа

\begin{equation}
\Delta U_l = 0.
\end{equation}

Визначимо тепер функцію $Y_l$ від відношень $x/r$, $y/r$, $z/r$ (тобто функцію, яка залежить лише від кутів $\theta$ і $\phi$) за допомогою рівняння

\begin{equation}
U_l = r^l Y_l.
\end{equation}

Далі підставимо цю функцію в рівняння Лапласа $\Delta U_l = 0$. Виконавши диференціювання за $r$, приходимо до рівняння

\begin{equation}
\left( \frac{\partial^2}{\partial r^2} + \frac{2}{r} \frac{\partial}{\partial r} \right) r^l Y_l + \frac{\Lambda}{r^2} r^l Y_l = r^{l-2} \left\{ \Lambda Y_l + l(l+1) Y_l \right\} = 0.
\end{equation}

Отже, введені вище функції є розв'язками диференціального рівняння $\Lambda Y_l + \lambda Y_l = 0$ лише за умови

\begin{equation}
\lambda = l(l+1).
\end{equation}

\begin{tcolorbox}[vintagebox, title={\upshape\rule{4mm}{0.3pt}~~Виведення власних значень оператора $\Lambda$~~\rule{4mm}{0.3pt}}]

Розглянемо задачу
\begin{equation}
\Lambda Y_l + \lambda Y_l = 0,
\end{equation}
де $\Lambda$~--- кутова частина лапласіана.
Лапласіан у сферичних координатах має вигляд
\begin{equation}
\Delta = \frac{\partial^2}{\partial r^2}
+ \frac{2}{r}\frac{\partial}{\partial r}
+ \frac{1}{r^2}\Lambda.
\end{equation}

\subsection*{Однорідні гармонічні багаточлени}
Нехай $U_l(x,y,z)$~--- однорідний багаточлен степеня $l$, що задовольняє рівнянню Лапласа
\begin{equation}
\Delta U_l = 0.
\end{equation}
Оскільки $U_l$ однорідний, його можна записати у вигляді
\begin{equation}
U_l(x,y,z) = r^l Y_l\!\left(\tfrac{x}{r},\tfrac{y}{r},\tfrac{z}{r}\right),
\end{equation}
де функція $Y_l$ залежить лише від кутів $\theta,\phi$.

\subsection*{Підстановка в лапласіан}
Підставимо $U_l = r^l Y_l$ у лапласіан:
\begin{equation}
\Delta U_l =
\left( \frac{\partial^2}{\partial r^2}
+ \frac{2}{r}\frac{\partial}{\partial r} \right) r^l Y_l
+ \frac{1}{r^2}\Lambda(r^l Y_l).
\end{equation}

Обчислимо радіальні похідні:
\begin{align}
\frac{\partial}{\partial r}(r^l Y_l) &= l r^{l-1} Y_l, \\
\frac{\partial^2}{\partial r^2}(r^l Y_l) &= l(l-1) r^{l-2} Y_l,
\end{align}
звідки
\begin{equation}
\left( \frac{\partial^2}{\partial r^2}
+ \frac{2}{r}\frac{\partial}{\partial r} \right) r^l Y_l
= l(l+1) r^{l-2} Y_l.
\end{equation}

Кутова частина дає
\begin{equation}
\frac{1}{r^2}\Lambda(r^l Y_l) = r^{l-2}\Lambda Y_l.
\end{equation}

Отже,
\begin{equation}
\Delta U_l = r^{l-2}\big( l(l+1) Y_l + \Lambda Y_l \big).
\end{equation}

\subsection*{Умова гармонічності}
Оскільки $U_l$ --- гармонічний багаточлен, то
\begin{equation}
\Delta U_l = 0,
\end{equation}
отже
\begin{equation}
\Lambda Y_l + l(l+1) Y_l = 0.
\end{equation}

\subsection*{Власні значення}
Ми отримали рівняння
\begin{equation}
\Lambda Y_l = -\,l(l+1) Y_l,
\end{equation}
тобто власні значення оператора $\Lambda$ дорівнюють
\begin{equation}
\lambda = l(l+1).
\end{equation}
Таким чином, функції $Y_l(\theta,\phi)$, що виникають з однорідних гармонічних багаточленів, є власними функціями кутового лапласіана, тобто сферичними (кульовими) гармоніками.

\end{tcolorbox}

Можна довести, що жодні інші значення $\lambda$ не дають скінченних, однозначних розв'язків цього рівняння. Відповідно до цього співвідношення
\begin{equation*}
    \Lambda Y_l + \lambda Y_l = 0
\end{equation*}
допустимі значення $\lambda$ дорівнюють $l(l+1)$. Крім того, число довільних параметрів у загальному розв'язку однозначно визначається степенем $l$. Наприклад, поліном степеня $l$ за $x$, $y$, $z$ містить $\frac{1}{2}(l+1)(l+2)$ довільних коефіцієнтів (він має один член з $x^l$, два члени з $x^{l-1}$, три з $x^{l-2}$ і т.\,д.; нарешті, $(l+1)$ членів взагалі не містять $x$). Очевидно, що умова $\Delta U_l = 0$ приводить до певної залежності між сталими; ця умова еквівалентна $\frac{1}{2} l(l-1)$ рівнянням для визначення коефіцієнтів, оскільки $\Delta U_l$ є однорідною функцією степеня $(l-2)$, тотожно рівною нулю. Тому $U_l$ містить

\begin{equation}
\frac{1}{2} \{ (l+1)(l+2) - l(l-1) \} = 2l + 1
\end{equation}
%
незалежних коефіцієнтів. Відповідно до цього існує $2l+1$ лінійно незалежних кульових функцій степеня $l$. Якщо записати їх у звичайній формі

\begin{equation}
Y_l^{(m)} = P_l^m e^{im\phi},
\end{equation}
%
то вони відповідають $2l+1$ можливим значенням третього (магнітного) квантового числа $m$.

\begin{note}
Іншими словами: при заданому орбітальному квантовому числі $l$ енергія електрона в кулонівському полі не залежить від того, як саме орієнтований його «орбітальний момент» --- значення $m$ може набувати $2l+1$ значень: $m = -l, -l+1, \dots, 0, \dots, l-1, l$. Це так зване \emph{виродження за $m$}: усі $2l+1$ станів з однаковими $n$ і $l$ мають однакову енергію (доки немає зовнішнього магнітного чи електричного поля).
\end{note}

Тепер перейдемо до диференціального рівняння для радіальної функції $R$:

\begin{equation}
\left\{ \frac{d^2}{dr^2} + \frac{2}{r} \frac{d}{dr} + \frac{8\pi^2 m}{h^2} \left( E + \frac{e^2 Z}{r} \right) - \frac{l(l+1)}{r^2} \right\} R = 0.
\end{equation}

Його розв'язки мають бути скінченними та неперервними при всіх значеннях $r$ від нуля до нескінченності. Нас цікавить, за яких значень енергії $E$ це рівняння взагалі має розв'язок, що задовольняє таким умовам --- адже для довільного $E$ такого розв'язку може не існувати. Розглянемо насамперед випадок $E<0$. У теорії Бора це відповідає еліптичним орбітам: щоб перевести електрон на нескінченну відстань від ядра, йому необхідно надати додаткову енергію (система зв'язана). Випадок $E>0$ відповідає гіперболічним (незв'язаним, розсіювальним) орбітам.

\begin{note}
Ключова ідея цього кроку методу: спочатку з'ясовуємо \emph{асимптотичну поведінку} розв'язку при $r\to\infty$ і при $r\to0$ окремо (це просто, бо там рівняння спрощується), а потім шукаємо повний розв'язок у вигляді добутку цих асимптотик на «повільну» невідому функцію $f(\rho)$. Це стандартний прийом розв'язання рівнянь типу Шрединґера з кулонівським потенціалом.
\end{note}

Для простоти введемо відносні (безрозмірні) одиниці. За одиницю радіуса приймемо борівський радіус $h^2/4\pi^2 me^2 Z$, а за одиницю енергії --- енергію основного стану атома Бора, $-2\pi^2 me^4 Z^2/h^2$. Іншими словами, покладемо (гл.~V, §~1)

\begin{equation}
    r = \rho \frac{h^2}{4\pi^2 me^2 Z}, \quad E = \varepsilon \left( \frac{-2\pi^2 me^4 Z^2}{h^2} \right).
\end{equation}

Тоді хвильове рівняння набуває простішого вигляду

\begin{equation}
\left\{ \frac{d^2}{d\rho^2} + \frac{2}{\rho} \frac{d}{d\rho} - \varepsilon + \frac{2}{\rho} - \frac{l(l+1)}{\rho^2} \right\} R = 0.
\end{equation}

Почнемо з визначення властивостей функції $R$ при дуже великих значеннях $\rho$. Відповідно відкинемо в диференціальному рівнянні члени з $1/\rho$ і $1/\rho^2$, так що поведінка на нескінченності визначатиметься рівнянням

\begin{equation}
\left\{ \frac{d^2}{d\rho^2} - \varepsilon \right\} R_\infty = 0.
\end{equation}

Воно має розв'язок

\begin{equation}
R_\infty = e^{\pm \rho \sqrt{\varepsilon}}.
\end{equation}

Проте розв'язок зі знаком «плюс» потрібно відкинути, оскільки в цьому випадку хвильова функція необмежено зростала б експоненційно при зростанні $r$ (або $\rho$) і тому не могла б бути власною функцією (не була б квадратично інтегровною).

Інша особлива точка знаходиться на початку координат. Насамперед, при дуже малих значеннях $r$ член $2/\rho$ прямує до нескінченності повільніше, ніж $l(l+1)/\rho^2$, тому в цій області справедливе наближене рівняння

\begin{equation}
\left\{ \frac{d^2}{d\rho^2} + \frac{2}{\rho} \frac{d}{d\rho} - \frac{l(l+1)}{\rho^2} \right\} R_0 = 0.
\end{equation}

Розв'язками його будуть функції $R_0 = \rho^l$ та $R_0 = \rho^{-l-1}$. Другий з них неприйнятний, оскільки перетворюється на нескінченність у нулі.

Отже, ми знаємо поведінку шуканої функції поблизу двох особливих точок, $\rho = 0$ і $\rho = \infty$. Розумно припустити, що вся функція $R$ має вигляд

\begin{equation}
R = e^{-\rho \sqrt{\varepsilon}} \rho^l f(\rho),
\end{equation}
%
де $f$ --- функція від $\rho$, яка, звісно, повинна бути регулярною в обох точках (тобто на нескінченності не повинна зростати швидше, ніж $e^{+\rho \sqrt{\varepsilon}}$, а в нулі --- залишатися скінченною). Множники $e^{-\rho\sqrt\varepsilon}$ та $\rho^l$ саме й забезпечують правильну асимптотичну поведінку на нескінченності та в нулі відповідно, тож уся «невідомість» переноситься на функцію $f(\rho)$.

Підставимо цей вираз для $R$ у диференціальне рівняння та прирівняємо коефіцієнти; отримуємо рівняння для $f$:

\begin{equation}
\frac{d^2 f}{d\rho^2} + \frac{2(l+1)}{\rho} \frac{df}{d\rho} - 2 \sqrt{\varepsilon} \frac{df}{d\rho} + \frac{2}{\rho} (1 - \sqrt{\varepsilon} (l+1)) f = 0.
\end{equation}

Спробуємо розв'язати його, розклавши $f$ у ряд за степенями $\rho$ (або, зручніше, за степенями $2\rho \sqrt{\varepsilon}$); запишемо відповідно

\begin{equation}
f = \sum_{v=0}^{\infty} a_v (2\rho \sqrt{\varepsilon})^v.
\end{equation}

Підставляючи цей вираз у диференціальне рівняння і дещо інакше групуючи члени, отримуємо

\begin{equation}
\sum_{v=0}^{\infty} a_v (2\rho \sqrt{\varepsilon})^{v-1} (v + 2l + 1) - \sum_{v=0}^{\infty} a_v (2\rho \sqrt{\varepsilon})^{v-1} \left( v + l + 1 - \frac{1}{\sqrt{\varepsilon}} \right) = 0.
\end{equation}

Цей ряд повинен тотожно перетворюватися на нуль; таким чином, ми приходимо до рекурентного співвідношення для коефіцієнтів

\begin{equation}
a_{v+1} (v + 1)(v + 2l + 2) = a_v \left( v + l + 1 - \frac{1}{\sqrt{\varepsilon}} \right).
\end{equation}

\begin{note}
Це рекурентне співвідношення --- ключ до всього результату. Воно виражає кожен наступний коефіцієнт ряду через попередній. Якщо ряд \emph{ніколи не обривається} (тобто містить нескінченно багато членів), можна показати, що при великих $v$ відношення $a_{v+1}/a_v \to 2\sqrt{\varepsilon}/v$, тобто ряд поводиться як розклад $e^{2\rho\sqrt{\varepsilon}}$. У такому разі $f(\rho)$ зростає при $\rho\to\infty$ швидше, ніж спадає множник $e^{-\rho\sqrt\varepsilon}$, і вся функція $R$ розбігається --- фізично неприйнятний (не нормовний) розв'язок. Єдиний спосіб уникнути цього --- щоб ряд \emph{обірвався}, перетворившись на скінченний поліном. Саме цю умову і накладено нижче.
\end{note}

Зрозуміло, що поблизу початку координат функція $f$ скінченна і дорівнює своєму першому члену. Проте, як показує детальніший аналіз, за великих $\rho$ функція $f$, що визначається нескінченним рядом, зростає швидше, ніж $e^{+\rho \sqrt{\varepsilon}}$ у всіх випадках, окрім того, коли ряд для $f$ обривається на якомусь члені, перетворюючись на скінченний поліном. Саме в цьому останньому випадку добуток $e^{-\rho\sqrt{\varepsilon}}\rho^l f(\rho)$, тобто $R$, на нескінченності прямує до нуля завдяки експоненційному множнику $e^{-\rho \sqrt{\varepsilon}}$. Умова, за якої відбувається обрив ряду, випливає з рекурентного співвідношення. Ряд обривається на $n_r$-му члені, якщо

\begin{equation}
n_r + l + 1 = \frac{1}{\sqrt{\varepsilon}}.
\end{equation}

Таким чином, $1/\sqrt{\varepsilon}$ повинно бути додатним цілим числом, тобто

\begin{equation}
\varepsilon = \frac{1}{n^2},
\end{equation}
%
де $n = n_r + l + 1$; число $n$ --- це головне квантове число, а $n_r$ --- радіальне квантове число.

\begin{note}
Отже, з умови «розв'язок має бути фізично прийнятним (скінченним усюди й нормовним)» \emph{автоматично} випливає, що енергія може набувати лише дискретного набору значень --- це і є квантування енергії, яке виникає тут природним чином із самої математики рівняння, без додаткових постулатів (на відміну від старої теорії Бора, де квантування вводилося як окремий постулат).

Зверніть увагу на співвідношення квантових чисел: оскільки $n_r = 0, 1, 2, \dots$ і $l = 0, 1, 2, \dots$, а $n = n_r+l+1$, то при заданому $n$ орбітальне число $l$ може набувати значень $l = 0, 1, \dots, n-1$ (тобто $n$ можливих значень), а кожному $l$ відповідає $2l+1$ значень магнітного числа $m$. Сумарне вироджування рівня $E_n$ дорівнює $\sum_{l=0}^{n-1}(2l+1) = n^2$ --- відомий результат для водневоподібного атома.
\end{note}

Отже, ми бачимо, що розв'язки диференціального рівняння існують лише для деяких дискретних значень параметра $\varepsilon$, а саме для значень $\varepsilon = 1/n^2$. Отже, можливі лише певні енергетичні рівні, а саме

\begin{equation}
E_n = - \frac{2\pi^2 me^4 Z^2}{h^2 n^2} = - \frac{h c R_\infty Z^2}{n^2},
\end{equation}
%
які й дає теорія Бора.

\begin{note}
Тут $R_\infty$ --- стала Рідберга (для нескінченно важкого ядра), $c$ --- швидкість світла; величина $hcR_\infty$ має розмірність енергії. Це саме той знаменитий результат, який Бор уперше отримав 1913 року з напівкласичної моделі й який чудово узгоджується з експериментальним спектром водню (формула Рідберга--Бальмера). Тут же він отриманий строго, як наслідок розв'язання рівняння Шрединґера.
\end{note}

Для $n=1$ маємо $l=0$, $n_r=0$, і $f$ зводиться до сталої. Хвильова функція $\psi$ не залежить від $\theta$ і $\phi$, і якщо прийняти нормування

\begin{equation}
\int \psi^2 d\tau = 1,
\end{equation}
%
де інтегрування ведеться по всьому простору ($d\tau = r^2\sin\theta\, dr\, d\theta\, d\varphi$), отримаємо просто

\begin{equation}
\psi = \frac{1}{\sqrt{\pi}} \left( \frac{Z}{a_1} \right)^{3/2} e^{-Z r / a_1}, \quad a_1 = \frac{h^2}{4 \pi^2 m e^2}.
\end{equation}

Можна побачити, що $\psi$ швидко зменшується при $r > a_1/Z$. Величина $a_1/Z$ --- це радіус першої борівської орбіти для атома із зарядом $Ze$.
%
Він дорівнює середньому значенню (див. додаток 25) $\overline{r} = \int \psi^2 r\, d\tau$ для електронної хмари в стані $n=1$. Такі середні величини можна обчислити для всіх вищих станів з $l=0$, і, як виявилося, вони збігаються з великими півосями еліпсів теорії Бора.

Можна додати, що поліноми $f$ --- це відомі поліноми Лагерра. Однак ми не будемо заглиблюватися в подробиці; зауважимо лише, що їхні нулі визначають розташування вузлових поверхонь. Функція $R$ має $n_r$ вузлів, не рахуючи нуля при $r=0$ (у випадку $l>0$) і при $r=\infty$.

\begin{note}
\textbf{Підсумок.} Розв'язання рівняння Шрединґера для кулонівського потенціалу дало три квантові числа:
\begin{itemize}
    \item $n = 1, 2, 3, \dots$ --- головне квантове число (визначає енергію $E_n \propto -1/n^2$);
    \item $l = 0, 1, \dots, n-1$ --- орбітальне (азимутальне) квантове число (визначає величину орбітального моменту імпульсу, $|\mathbf{L}|=\hbar\sqrt{l(l+1)}$);
    \item $m = -l, \dots, l$ --- магнітне квантове число (визначає проєкцію моменту імпульсу на вибрану вісь, $L_z = m\hbar$).
\end{itemize}
Радіальна частина хвильової функції виражається через поліноми Лагерра, кутова --- через сферичні функції $Y_l^{(m)}(\theta,\varphi)$, побудовані з приєднаних поліномів Лежандра.
\end{note}

\end{document}
